\par\bigskip\Needspace{7\baselineskip}
\pdfbookmark[1]{Randbemerkungen in Gauss' Handexemplar}{abbildung-randbemerkungen}
\begin{center}[\textit{Randbemerkungen in} \textsc{Gauss}' \textit{Handexemplar}:]\end{center}
\begingroup\small
\par\indent[Artikel 10, neben der letzten Gleichung zur Bestimmung von \(\lambda\)] oder \(\lambda=\cos u^*\), wenn der kleinste Wert [des Vergrößerungsverhältnisses] bei \(u=u^*\) angenommen werden soll.

[Artikel 12, neben der Gleichung, mit der in diesem Druck die Größe \(w\) eingeführt wird] Das Zeichen \(\omega\) ist entgegen meiner Absicht gedruckt worden. Es sollte \(w\) sein.

[Artikel 13, neben den Gleichungen zur Abbildung, die durch die Funktion \(f\upsilon=\upsilon-i\log k\) bestimmt ist, stehen die entsprechenden Gleichungen für die Funktion \(f\upsilon=a\upsilon-i\log k\). Sie wurden später in die erste Abhandlung der \textit{Untersuchungen über Gegenstände der höhern Geodäsie} aufgenommen.]
\par\endgroup
\begin{center}\rule{1.2in}{.3pt}\end{center}
